\chapter{The role of the writer is not simply to arrange (There's No Way I'll Ever Regret It)}
In Episode 5, the beginning of the second arc continues to mirror the first arc. Just as Episodes 1 and 4 both served as introductions, Episodes 2 and 5 are both about positioning the characters and exploring the nature of this new world, which is a polite way of saying that this is an episode where not much happens.

The episode opens with a flashback to Kyubey and Sayaka performing the ritual that transforms her into a magical girl, presumably right after she assured Kyousuke that magic exists in Episode 4. While Kyubey has been creepy throughout the series so far, this flashback is the longest sustained depiction thus far of him as a (literally, here) dark figure, and the framing and lighting both are highly suggestive of a death scene. As Homura states later, Sayaka's fate is fixed at this point; she is, effectively, a dead woman walking. (Again, quite literally, as we will learn in a few episodes.) Further, by placing Kyubey in deep shadow with a large plant behind him, several shots look as if Kyubey has multiple tails, visually resembling the powerful, nine-tailed spirit fox, a (usually malevolent) trickster spirit \cite{ref40}. 

We also see for the first time how a Soul Gem is formed. The ritual suggests strongly that Kyubey literally pulls it out of the girls' hearts, making it from something that already existed within them. This accords with statements by Mami and Kyubey in prior episodes that they can sense great power in Madoka, even though she has yet to become a magical girl, and explains why Kyubey does not simply use his power to accomplish his goals: though enormous, his power is extremely limited. He can grant wishes, act as a telepathic switchboard, control who can see him, and (as we will learn later in the series) exist in multiple places at once, but cannot actually wield magic to alter reality the way the magical girls and witches do. A few of his comments even suggest that what wishes he can grant is determined by the power of the magical girl doing the wishing—given comments in later episodes that Sayaka is not a very powerful magical girl, it's possible that the reason she only wished for Kyousuke's hand to be healed and not the rest of his body is that she couldn't heal the rest.

Kyubey, in other words, is a facilitator. He enables prospective magical girls to tap a power that already exists in them, so that they can fight witches or each other for him. As we see in this episode, he is perfectly happy to construct a conflict, empowering Sayaka even though he knows Kyoko is coming, feeding Kyoko information while keeping Sayaka in the dark, all because the fight between them serves his ends. 

As a consequence of Kyubey's manipulations, Kyoko takes over Homura's role in the first arc as the antagonistic magical girl of questionable morality. Kyoko is everything Mami warned about: highly willing to fight other magical girls, concerned only with the rewards of defeating witches, and uncaring about protecting the people of Mitakihara. Her willingness to let the familiar kill people until it becomes a witch, along with her comments regarding the food chain and her own constant eating, combine to suggest that Kyoko sees eating as an expression of power and embraces a might-makes-right philosophy regarding that power. In opposing her, Sayaka takes over Mami's role as the "good" magical girl, the one who fights to protect others and believes the strong have a duty to defend the weak.

By interrupting them, Homura reveals that she has taken over the other role Mami played in the first arc. Homura is no longer trying to erase the traditional magical girl structure and replace it with the show Madoka Magica will be; that has already been achieved. Instead, she is now trying to prevent the next logical development in the story, the Magical Girl Madoka promised by the title. She refuses to help Sayaka when Madoka begs her to do so, but when the only alternative is for Madoka to become a magical girl, Homura has no choice but to step in.

This leaves only Madoka and Kyubey. Madoka makes an interesting and deliberate choice to not change her role—just as she was Mami's unpowered sidekick and confidante, so she offers to be for Sayaka. Even if she is understandably terrified of becoming a magical girl, she is still willing to risk her life to stand by Sayaka's side—and as we see at the end of the episode, even willing to become a magical girl if need be. Kyubey likewise does not change his role, but rather increasingly reveals to the audience what his role is, moving from a figure of questionable morality and allegiance to an obviously manipulative figure who is increasingly positioned as antagonistic, actively assisting Kyoko and keeping her a secret from Madoka and Sayaka.

But what precisely is that role? Early in the episode, Sayaka takes Kyousuke to the roof, where his family give him back his violin and he plays "Ave Maria." What makes this notable is that elements of the scene keep comparing Kyousuke to Kyubey. First, this is the same location in which Sayaka made the pact to become a magical girl at the beginning of the episode, and although they are outside rather than inside the concentric rings of flowers, Sayaka and Kyousuke are in the same relative positions as Sayaka and Kyubey were in that flashback. Second, Kyousuke is shown as a silhouette in some shots, just as Kyubey is in the flashback, and in shots from Sayaka's point of view, Kyousuke blocks the view of where Kyubey was standing in the flashback. Most subtly, but also most importantly for understanding the function of the Kyousuke-Kyubey parallel, Kyubey reaches into Sayaka's heart to make her into a magical girl, just as Kyousuke's music reaches into Sayaka's heart, creating her feelings for him that motivated her to become a magical girl. And, of course, we will see in the next couple of episodes that Kyousuke is rather thoughtless in his behavior toward Sayaka, so he shares with Kyubey that they facilitated Sayaka's transformation into a magical girl while caring very little about her feelings.

Near the end of the scene, Kyousuke puts down the violin, and yet "Ave Maria" continues to play as the background music for the rest of the scene, transitioning from diegetic (that is, "in-universe") to extradiegetic sound. This ability to straddle the borders of diegesis has, up until this point in the show, been presented as an ability possessed by the witches. To convey their otherworldliness, the witch's labyrinths are generally given their own unique art styles, distinct from the show, with the result that we see the characters noticing, and reacting with terror to, changes in the art style—a diegetic response to an extradiegetic effect. Kyousuke is now reaching across that barrier in the opposite direction, one of only three non-witch characters to cross that threshold—and the only one without any apparent "magical" abilities.

This is because his music is an expression of emotion, and therefore magic; as we will see much later in the series, human emotion is the source of all the magic in the series. It follows that human art is therefore fundamentally magical; that artistic expression can reshape reality. Certainly it has done so here: Kyousuke's music is heavily implied to be the source of Sayaka's interest in him, which is the reason she became a magical girl. Every violation of the laws of physics performed by Magical Girl Sayaka is thus a consequence of Kyousuke's music.

But if Kyousuke can cross between diegetic and extradiegetic in a scene in which he is heavily paralleled with Kyubey, does it follow that Kyubey can do likewise? Indeed, he can, and is the second of the three non-witch characters to do so. Kyubey is an extradiegetic entity taking up residence in the story.

Consider: Kyubey creates magical girls to serve his own purposes, knowing that they will suffer—even relying on that suffering. He wants Madoka to become a magical girl, and shapes everything he does toward that end result, since he has a problem she can help solve. He sets up Kyoko to fight Sayaka for the same reason, once again caring nothing for how it effects them except insofar as those effects serve his goals. What are those goals? To bring tragedy and despair to the magical girls he creates.

In this sense he is heavily reminiscent of the villain of another magical girl series, Drosselmeyer of Princess Tutu \cite{ref41}. In that series, Drosselmeyer is ultimately revealed to be an author with the power to bring his creations to life, which he uses to trap a town in a cycle of tragedy \cite{ref42}. However, the pure sacrifice of the duck he transformed into the title character brings about the end of cycle of tragedy and forces a happier ending, sending Drosselmeyer off into other wor(l)ds. Compare this plot to that same afterword of Urobuchi's we discussed last chapter: "I am full of hatred towards men's so-called happiness, and had to push the characters I poured my heart out to create into the abyss of tragedy." And later, "Only a heavenly and chaste soul that can sing carols of praise towards humanity can save the story." \cite{ref43} (And see also Chapters 20 and 21 for more relationships between Kyubey and the probable source for Princess Tutu's Drosselmeyer, E.T.A. Hoffman's character Drosselmeier.)

Among other things, Madoka Magica is a series about consent and autonomy. There have been hints toward this theme, but it becomes undeniable in the next episode. Given that, what better villain for such a series than the one who controls the actions of the characters? All of Kyubey's abilities align: to make magical girls, to know what they're thinking, to be everywhere in the story at once. All of his motivations—to make the magical girls experience emotional highs and lows, to keep the world of the story running as long as possible—are likewise consistent with the role he plays without knowing it: Kyubey is the author of Puella Magi Madoka Magica.
